\documentclass[preprint,12pt,numbers]{elsarticle}
\usepackage{amsmath,amssymb,amsthm}
\usepackage[hidelinks]{hyperref}
\newtheorem{theorem}{Theorem}[section]
\newtheorem{lemma}[theorem]{Lemma}
\newtheorem{proposition}[theorem]{Proposition}
\newtheorem{corollary}[theorem]{Corollary}
\makeatletter
\let\ps@pprintTitle\ps@plain
\makeatother
\begin{document}
\begin{frontmatter}
\title{Blocking vertical translates of monomial graphs and cuspidal curves}
\author{GPT-6.2 Astra}
\begin{abstract}
We determine the order of the minimum number of vertical translates of
the graph $y=x^d$ needed to meet all rational projective lines over
$\mathbb F_q$. For each fixed $d\ge2$ this number is
$\Theta_d(\log q)$ uniformly in the characteristic. The same order holds
for the rational curves parametrized by $(t^{d-1},t^d)$. For cubic
permutations in characteristic greater than three, a two-parameter
count gives a lower limit of at least $1/\log 4$ after division by
$\log q$, for both pencils. The two minima can nevertheless differ
by an unbounded factor on the same field and in the same degree.
Over $\mathbb F_{s^2}$ in degree $s$ the cuspidal pencil needs three
translates and the graph needs $2s-1$ for every prime power $s$.
Over $\mathbb F_{s^3}$ the corresponding exact values are four and
$3s-2$ for $s\ge4$. These comparisons follow from an obstruction
involving translates of norm fibres.
\end{abstract}
\begin{keyword}
finite field \sep blocking set \sep monodromy \sep character sum \sep norm equation
\MSC[2020] 11T06 \sep 11T23 \sep 14G15 \sep 51E21
\end{keyword}
\end{frontmatter}

\section{Two pencils of rational curves}
Let $F=\mathbb F_q$, and let $d\ge2$. A subset of the rational
projective plane blocks if it meets every rational line. We consider
the two pencils
\begin{align*}
 G_\alpha:\quad YZ^{d-1}&=X^d-\alpha Z^d,\\
 C_\alpha:\quad X^d&=(Y+\alpha Z)^{d-1}Z,\qquad \alpha\in F .
\end{align*}
Write $B(q,d)$ and $B_{\rm cusp}(q,d)$ for the least numbers of parameters
whose corresponding unions block. In degree two the members of the
second pencil are smooth conics. Both pencils have the single point
$[0:1:0]$ at infinity, so all vertical lines and the line at infinity
meet every member. All logarithms are natural.

The graph minimum is Problem~5.4 of Asgarli, Ghioca and Yip \cite{AGY}.
Their Theorem~5.1 gives its logarithmic upper bound, and their
Remark~5.3 gives a logarithmic lower bound when $\gcd(d,q-1)>1$.
Indeed, the horizontal lines already require translates of the
set of $d$th powers to cover $F$. Bounds for additive translates of
multiplicative subgroups are developed further in \cite{ZhouWu}.
In permutation degree, a single member meets every horizontal line.
Theorem~\ref{thm:fixed-graph} obtains a logarithmic lower bound in
this case as well, even if all translated covers share their only
branch value. The question about arbitrary collections of plane curves
in \cite[Thm.~1.3]{AGY} has different quantifiers.

The second pencil has a separate source of monodromy. Over every
characteristic, one finite fibre of $t^d-mt^{d-1}$, $m\ne0$, has
multiplicities $d-1$ and $1$. For $d\ge3$ the inertia at this fibre
gives a doubly transitive subgroup supported in a single coordinate
after separation of the translates. The resulting argument for the
cuspidal pencil uses no classification of finite doubly transitive
groups. For both pencils the lower limit after division by $\log q$
is at least $1/(2\log(d!))$. For fixed $d$ the blocking sets obtained
with $O_d(\log q)$ translates are nontrivial for large $q$: each
line meets each geometrically irreducible member in at most $d$ points.

For $d=3$, $q\equiv2\bmod 3$ and $\operatorname{char}F>3$, the
two-parameter estimate in Theorem~\ref{thm:cubic-surface} gives the
stronger coefficient $1/\log4$. The estimate is uniform in the
translate set, including its cardinality. The final section lets the
degree vary and computes both minima on the same field. In quadratic
and cubic extensions the cuspidal obstruction has a description by
norm equations that gives a sharp finite result independently of the
point-counting arguments.

\subsection{A uniform point count on a twist}
\begin{lemma}\label{lem:twist}
Let $P_i\in F[x]$, $1\le i\le\ell$, be separable as maps of degree
$d\ge2$ (so $P'_i$ is not identically zero), and let $M$ be the
splitting field of the product of $P_i(x)-t$ over $F(t)$.
Let $H$ be its geometric group and $b$ the number of geometric
branch values, including infinity. Suppose the arithmetic group
contains an element $\gamma$ acting on the full constant field as
$q$-Frobenius and fixing no root in any coordinate. If
\begin{equation}\label{eq:twist-budget}
 q+1-\ell(d-1)|H|\sqrt q>b|H|,
\end{equation}
some $t_0\in F$ has the property that no $P_i(x)=t_0$ has a root in $F$.
\end{lemma}
\begin{proof}
Over $k=\overline F$, adjoin the roots of each polynomial one at a
time. A single root $x$ generates a rational field with
$[k(x):k(t)]=d$. If the present field $K$ has degree $N$ over $k(t)$,
let $K'=Kk(x)$, $e=[K':K]$ and $N'=eN$. Then
$[K':k(x)]=N'/d$. The Castelnuovo--Severi inequality
\cite[Thm.~3.11.3]{Stichtenoth} gives
\[
 \frac{g(K')-1}{N'}\le\frac{g(K)-1}{N}+\frac{e-1}{d}.
\]
The $j$th root of one polynomial costs $e\le d-j+1$, including
steps of degree one. Indeed, the $j-1$ previously adjoined
distinct roots leave a polynomial of degree $d-j+1$.
With $N=1$ and $g=0$ at the start, the sum telescopes to
\[
        2g+2|H|\le\ell(d-1)|H|+2 .
\]
Dropping $2(|H|-1)$ gives the bound in \eqref{eq:twist-budget}.
The argument counts the full genus, including any wild ramification at infinity.

Write $F'$ for the constants of $M$ and $a$ for the order of
$\gamma$. Then $F'\subseteq\mathbb F_{q^a}$. On $M\mathbb F_{q^a}$
combine $\gamma$ with $q$-Frobenius on $\mathbb F_{q^a}$, which
agree on $F'$. The fixed field of the combined automorphism gives
the twist over $F$ \cite[Thm.~2]{Kosters}. Its full constant field
is $F$ and base change to $k$ recovers the geometric splitting
field. Its smooth projective curve has genus $g$ and maps to the
line with degree $|H|$. By Hasse--Weil it has at least
$q+1-2g\sqrt q$ rational points. At most $b|H|$ geometric points lie
over the excluded values. A remaining rational point specializes
to the Frobenius class of $\gamma$ (or its inverse with the opposite
convention); its cycles are the factor degrees in each fibre. This
proves the lemma.
\end{proof}

For reference, both pencils have $|H|\le D^\ell$, $D=d!$, and
$b\le1+d\ell$. A finite sufficient window for \eqref{eq:twist-budget}
is $L=\log q\ge1$ and
\begin{equation}\label{eq:fixed-window}
       \ell\log D\le L/2-\log L-\log(4(d+1)).
\end{equation}
Indeed $\ell\le L$ since $\log D\ge\log2>1/2$, and
$4(d+1)L D^\ell\le\sqrt q$. The genus loss is at most $q/4$ and
the discarded fibres cost at most $\sqrt q/4$. The separate slope
exclusions below are polynomial in $\ell$ for fixed $d$. In
particular their validity is uniform throughout this window for
all sufficiently large $q$.

\section{Graph pencils in every characteristic}\label{sec:fixed-graph}

\begin{theorem}\label{thm:fixed-graph}
For every fixed $d\ge2$, along all prime powers $q$,
\[
 B(q,d)=\Theta_d(\log q),\qquad
 \liminf_{q\to\infty}\frac{B(q,d)}{\log q}\ge \frac1{2\log(d!)}.
\]
\end{theorem}

The upper bound is Theorem~5.1 of Asgarli--Ghioca--Yip
\cite{AGY}. Their Remark~5.3 obtains a logarithmic lower bound from the
horizontal lines when $\gcd(d,q-1)>1$. We treat all characteristics and
both permutation and nonpermutation degrees.

First consider $d=p^a$ and $q=p^h$. Put $g=\gcd(a,h)$, $s=p^g$ and
$n=h/g$. The automorphism $\sigma(x)=x^d$ generates the Galois group
of $\mathbb F_q/\mathbb F_s$. For each slope $m$ the map
$\sigma-m\,\mathrm{id}$ is invertible or has one-dimensional kernel.
In the latter case $m=\sigma(u)/u$ for some $u\ne0$, and its image is
$\sigma(u)\operatorname{im}(\sigma-1)$. The latter unscaled image is
the trace-zero hyperplane by dimension and the trace identity. By the
nondegenerate trace pairing its scalar multiples range over all linear
hyperplanes. Hence a set of translates blocks precisely when it meets
every affine $\mathbb F_s$-hyperplane. The Jamison--Brouwer--Schrijver
theorem gives at least $n(s-1)+1$ points \cite{Jamison,BS}; the union
of the $n$ coordinate axes meets the bound. Consequently
\begin{equation}\label{eq:graph-frob}
                 B(p^h,p^a)=1+\frac hg(p^g-1).
\end{equation}
Its logarithmic coefficient $(s-1)/\log s$ is at least one.

For the remaining degrees set $k=\overline{\mathbb F}_q$ and
$f_m(x)=x^d-mx$. A missing line for $S$ is an intercept $b$ and a slope
$m$ for which no $f_m(x)=b+\alpha$ has a rational root. We construct
the required element of a Frobenius coset before applying
Lemma~\ref{lem:twist}. The assertions about ramification in this
section refer to geometric points.

\paragraph{Separated finite branch loci.}
Suppose $p\nmid d(d-1)$, which forces $p\ne2$. For
$m=d u^{d-1}$ the finite critical values are $(1-d)u^d\zeta$ for
$\zeta^{d-1}=1$. Their simple critical points give transposition inertia
and infinity is tame of index $d$. For each distinct pair of shifts and
each ordered pair of distinct roots of unity, a collision excludes
at most $d$ choices of $u\in\mathbb F_q^\times$. Thus, once
$q-1>d(d-1)(d-2)\binom{\ell}{2}$, all translated finite branch loci
can be separated. This includes $d=2$ in odd characteristic, for which
there is just one finite branch value per cover.

A connected cover of the projective line tame everywhere and unramified
away from infinity has degree one by Riemann--Hurwitz. The tame local
extensions are closed under conjugates and composita, including when
$p\mid d!$. The finite inertia groups therefore generate the single
geometric group normally. They are transpositions and their conjugates
generate a transitive group, hence $S_d$. The intersection of a new
splitting field with the preceding compositum is unramified at every
finite place and tame at infinity. It is trivial by the same argument.
The geometric group of the compositum is $S_d^\ell$. This use of
separated branch values for translated splitting fields belongs to
earlier work on polynomial correlations \cite{Kurlberg,Pollack}.

We shall use double transitivity in the other cases. Write $d=p^a v$
with $p\nmid v$, allowing $a=0$. In the cases still under consideration
$v>1$. For $f=x^d-x$ the geometric group $G$ is transitive, since
$k(x)/k(f(x))$ is separable of degree $d$. The distinct ordered root
pairs form a single orbit: their curve $f(x)=f(y)$, off the diagonal,
has on the dense open with $z=x/y$ the equation
\[
                 y^{d-1}=(z-1)/(z^d-1).
\]
At a nonidentity $v$th root of unity the right side has valuation
$-p^a$, prime to $d-1$. On adjoining any root $y$ the ramification
index there is consequently divisible by $d-1$, so the equation is
irreducible over $k(z)$. This valuation argument also works when the
equation in $y$ is inseparable. The excluded loci $xy=0$ contribute
no generic pairs. The resulting geometric irreducibility gives double
transitivity of $G$.

If $p\mid d-1$, all finite critical points of $f$ have value zero and
infinity is tame. Take $m=1$. The translated branch values $-\alpha$
are automatically distinct. The intersection argument proves that the
composite geometric group is $G^\ell$ even though the finite inertia
can be wild.

\paragraph{Only the common wild place.}
Now $p\mid d$ and $d$ is not a power of $p$. The derivative of $f$ is
$-1$. Write $L/k(t)$ for its splitting field. Let $E$ be the additive
group of $c\in k$ whose translation $t\mapsto t+c$ preserves $L$
inside an algebraic closure; normality makes this independent of the
extension of the translation. It acts on $G$-conjugacy classes of
intermediate fields. The class of $k(x)$, where $t=f(x)$, has trivial
stabilizer. Indeed, after composition with an element of $G$ a
stabilizing translation lifts to an automorphism of $k(x)$ fixing the
unique pole of $f$. It has the form $x\mapsto A x+B$, and
\[
             (Ax+B)^d-(Ax+B)=x^d-x+c .
\]
Here $a\ge1$, $v>1$ and $2\le d-p^a<d$. Comparison at $x^{d-p^a}$
gives $vA^{d-p^a}B^{p^a}=0$, whence $B=0$. The linear and constant
terms give $A=1$ and $c=0$. There are at most $2^{|G|}$ subgroups of
$G$, so $|E|\le2^{d!}$.

Choose $m=u^{d-1}$, write $x=uy$ and $b=u^dt$. Two normalized
splitting fields coincide precisely when $(\alpha-\beta)/u^d\in E$.
At most $d\,2^{d!}\binom{\ell}{2}$ nonzero values of $u$ are excluded.
All fields can thus be made pairwise distinct when $\ell=O_d(\log q)$
and $q$ is sufficiently large.

Since $d$ has at least two distinct prime divisors, the double
transitivity of $G$ excludes affine type. By the classification of
finite doubly transitive groups, $T\le G\le\operatorname{Aut}(T)$ with
$T$ nonabelian simple. Its socle $T$ is doubly transitive in this action
except possibly for $\mathrm{P}\Gamma\mathrm L_2(8)$ in degree $28$
\cite[\S2.2]{Reid}. In that exception $G/T$ has order $3$ and
$p\in\{2,7\}$. This gives a cubic quotient cover unramified away from
infinity and tame there, contradicting Riemann--Hurwitz. Thus $T$
is doubly transitive.

The composite geometric group $H$ is a subdirect subgroup of $G^\ell$.
Its last derived subgroup $P$ lies in $T^\ell$ and projects onto $T$
in each factor. Indeed, the last derived subgroup commutes with
surjective images, $T$ is perfect, and $G/T$ is solvable by the
Schreier theorem. If a two-coordinate projection of $P$ is proper,
Goursat's lemma makes it the graph of an automorphism of $T$.
Normality of $P$ in $H$ then implies that for each $h\in H$ its two
coordinate conjugation actions on $T$ correspond under that
automorphism. Since $C_G(T)=1$, one coordinate of $h$ is the identity
exactly when the other is. The two kernel fixed fields coincide,
contrary to the choice of $u$.

It follows that $P=T^\ell$. To see this, project onto the first
$\ell-1$ coordinates and use induction. The kernel in the last factor
is normal in $T$. A trivial kernel would make the last projection
a surjection $T^{\ell-1}\to T$. Its commuting normal factor images
force it to factor through one coordinate, contradicting full pair
projections. The other possibility is the full kernel $T$. We have
proved $H\supseteq T^\ell$ without assuming independence of the
solvable quotients.

\paragraph{Frobenius and uniformity.}
In every separable case the geometric group contains $J^\ell$ for a
doubly transitive group $J$ on the $d$ roots. In the arithmetic group
choose $\sigma$ inducing rational Frobenius on the full constant field.
For any permutation $\tau$ of the roots and uniform $j\in J$, the
number $Z$ of fixed points of $\tau j$ has
$\mathbb E Z=\mathbb E[Z(Z-1)]=1$, by transitivity on points and on
ordered distinct pairs. If $Z\ge1$ always, it would equal one always,
contradicting the second identity. We may therefore choose the
coordinates of $j\in J^\ell$ independently so that every coordinate of
$\sigma j$ is a derangement.

There are at most $1+(d-1)\ell$ geometric branch values, including
infinity. With $D=d!$, Lemma~\ref{lem:twist} bounds the ratio of the
point-count error to $q$ by
\[
    (d-1)\ell D^\ell/\sqrt q+
    \big(1+(d-1)\ell\big)D^\ell/q .
\]
For $\ell=\lfloor c\log q\rfloor$ with $c<1/(2\log D)$, both terms
tend to zero uniformly in $S$ and $p$, and the exclusions on $u$ are
$O_d(\ell^2)=o(q)$. An unramified rational intercept with the chosen
Frobenius class exists, giving a missing line. This proves
Theorem~\ref{thm:fixed-graph}. The inequality is obtained from the
point count on the twist, so possible extensions of the constants have
already been included.

\section{The pencil with adjacent exponents}\label{sec:cusp-general}

Let $B_{\rm cusp}(q,d)$ denote the minimum number of members of
\[
                 X^d=(Y+\alpha Z)^{d-1}Z,\qquad \alpha\in\mathbb F_q,
\]
whose rational points meet every rational projective line. For $d=2$
these are conics. For $d>2$ the affine origin of the untranslated
member is a cusp. Each form is geometrically irreducible: in the
coordinates $X,V=Y+\alpha Z,Z$, it is primitive and linear in $Z$
over $\overline{\mathbb F}_q[X,V]$.

\begin{theorem}\label{thm:cusp-general}
For each $d\ge2$, uniformly over all characteristics,
\[
 B_{\rm cusp}(q,d)=\Theta_d(\log q),\qquad
 \liminf_{q\to\infty}\frac{B_{\rm cusp}(q,d)}{\log q}
                 \ge\frac{1}{2\log(d!)} .
\]
\end{theorem}

\begin{proof}
The affine points are $(t^{d-1},t^d-\alpha)$, including the point
at $t=0$: for $x\ne0$ the inverse parameter is $(y+\alpha)/x$.
The unique point at infinity is $[0:1:0]$. Thus the translates
indexed by $S$ block precisely when for every slope $a$ and intercept
$b$ the set $b+S$ meets the image of $f_a(t)=t^{d-1}(t-a)$.
Every image has at least $q/d$ elements, since every fibre has size
at most $d$. For $k=\lceil2d\log(2q)\rceil$ independent uniform
shifts, the expected number of missed lines is at most
$q^2(1-1/d)^k\le1/4$. Removing repeated shifts proves the upper
bound by the same value-set covering principle as in
\cite[Theorem~5.1]{AGY}. For large $q$ the resulting blocking set is nontrivial:
each line contains at most $dk$ points of the union by B\'ezout.

For $d=2$ this pencil equals the quadratic graph pencil, so the
lower bound follows from the quadratic case. Suppose $d\ge3$ and
fix a set $S$ of $\ell$ distinct shifts. Take $a\ne0$. The derivative
of $f_a$ is $t^{d-2}(dt-(d-1)a)$, a nonzero polynomial in every
characteristic. In the zero fibre the points $t=0,a$ have
ramification indices $d-1,1$. There is at most one other finite
branch value. It occurs only when $d(d-1)$ is nonzero in the field,
and then equals $\kappa a^d$ with
$\kappa=-(d-1)^{d-1}/d^d\ne0$. The finite branch values of
$b=f_a(t)-\alpha$ are these values translated by $-\alpha$.
If a second value occurs, exclude $\kappa a^d=\alpha-\beta$ for
ordered distinct $\alpha,\beta\in S$. This excludes at most
$d\ell(\ell-1)$ slopes. Whenever $q-1>d\ell(\ell-1)$ some slope
remains, and each $-\alpha$ belongs to the branch locus of just
one factor.

Work first over $\overline{\mathbb F}_q(b)$. Let $H$ be the Galois
group of the compositum of the splitting fields of
$f_a(T)-b-\alpha$ and $G_\alpha$ its projection on the roots of one
factor. This action is transitive, since the rational map $f_a$ is
separable of degree $d$. Fix a place above $b=-\alpha$ in the joint
splitting field. Its inertia acts trivially on the residue field
\cite[Tag 09EC]{Stacks}. All roots are integral. In every other
coordinate their residues are distinct by branch separation, so
inertia fixes each root. In the chosen coordinate a single root
reduces to $a$ and the other $d-1$ reduce to zero.

Here is a local construction of the required permutations on the latter
roots. Put $m=d-1$, $h=a+u$ and $b+\alpha=u h^m$. This is a uniformizing
parameter in $k[[u]]$, since $a\ne0$. The chosen polynomial factors as
\[
  T^m(T-a)-uh^m=(T-h)Q_m(T),\qquad
  Q_0=1,\quad Q_{j+1}=TQ_j+uh^j .
\]
The reduction of $Q_m$ is $T^m$ and its constant coefficient
$uh^{m-1}$ has order one. Hence $Q_m$ is Eisenstein over $k[[u]]$.
Its linear coefficient $uh^{m-2}$ is nonzero for $m\ge2$, so the
irreducible factor is separable over $k((u))$, also when $p\mid m$.
Every other coordinate splits over $k((u))$ by Hensel lifting of
its distinct residue roots over $k$. A local automorphism can
therefore carry any root of $Q_m$ to any other while fixing $h$
and every other coordinate. Restrict this automorphism to the normal
global splitting field embedded in the same local algebraic closure.
In the usual decomposition-group description the local Galois group
is inertia because the constants are algebraically closed. The
restriction construction itself already gives the required global
permutations.

Consequently the normal closure of the local image in $H$ is supported
in coordinate $\alpha$. Its projection $N_\alpha$ is normal in
$G_\alpha$ and doubly transitive. Indeed, conjugation carries the
inertia to a subgroup fixing any specified letter and transitive
on its complement. To move one letter to another, fix a third
letter; then a point stabilizer permits any prescribed move on
the remaining letters. In particular $H$ contains the product of
these $N_\alpha$ in their separate coordinates.

Let $\sigma$ be a lift of the $q$-Frobenius on the full constant
field of the arithmetic compositum. In every coset
$\sigma_\alpha N_\alpha$ the average number of fixed letters is
one by transitivity: each letter is carried to each given target
equally often. Double transitivity supplies a member fixing two
prescribed letters, so another fixes none. Multiplying $\sigma$ by
an element of $\prod_\alpha N_\alpha$ gives an element $\gamma$ in
the required arithmetic coset with no fixed letter in any coordinate.

Put $D=d!$. Each root $x$ generates a rational field of degree $d$
over the base, and $|H|\le D^\ell$. Lemma~\ref{lem:twist} therefore
applies to $\gamma$, including shared wild ramification at infinity
and the possibility of extra constants. There are at most $2\ell+1$
geometric branch values including infinity. The twist has an
unramified rational point, giving a missing line, as soon as
\[
 q+1>\ell(d-1)D^\ell\sqrt q+(2\ell+1)D^\ell.
\]
For $\ell\le c\log q$ with fixed $0<c<1/(2\log D)$ the two terms
on the right, divided by $q$, are
$O_d((1+\ell)q^{-\epsilon})$ and
$O((1+\ell)q^{-1/2-\epsilon})$ respectively, where
$\epsilon=1/2-c\log D>0$. The slope exclusion is $O_d(\ell^2)$.
All three estimates are independent of the characteristic and of
$S$, as required.
\end{proof}

\section{Two moving parameters for cubics}\label{sec:cubic-surface}

Throughout this section $p>3$ and $q\equiv2\pmod3$. For a nonempty
set $S\subset\mathbb F_q$ of cardinality $\ell$, let
\[
 X=\mathbb A^1\setminus(-S),\qquad U=(X\times X)\setminus\Delta,\qquad
 t_\alpha=\frac{2(v+w+2\alpha)}{v-w}.
\]
Write $N(S)$ for the number of $(v,w)\in U(\mathbb F_q)$ for which
$T^3-3T-t_\alpha$ is irreducible for every $\alpha\in S$.

\begin{theorem}\label{thm:cubic-surface}
Uniformly for all $q$ and $S$ as above,
\begin{align}
 \big|3^\ell N(S)-(q-\ell)(q-\ell-1)\big|
 &\le \ell(\ell-1)(4^\ell-1)q \notag\\
 &\quad +10\ell\,2^\ell(q+q^{3/2}).\label{eq:cubic-error}
\end{align}
With $L=\log q$, both the graph and the cusp blocking numbers in degree
three are greater than $(L-2\log L)/\log4-5$ for all sufficiently
large $q$. For each family the upper limit after division by $L$ is
at most $2/\log3$.
\end{theorem}

For the graph, let $4u^3=v-w$, which has a unique solution, and put
$m=3u^2$, $b=(v+w)/2$. Then $x=uT$ converts the intersection equation
$x^3-mx=b+\alpha$ into $T^3-3T=t_\alpha$. Thus a pair counted by $N(S)$
gives a missing line. Each missing line with $m\in3(\mathbb F_q^\times)^2$
arises from the two orders of its pair.

For the cusp, the affine parametrization is $(t^2,t^3-\alpha)$. This time
set $m=3u$, $w=b$, $v=b+4u^3$. All nonzero slopes and intercepts occur
exactly once. Substituting $t=u(T+1)$ gives
\[
  t^3-mt^2-b-\alpha =u^3(T^3-3T-t_\alpha).
\]
If $w=-\alpha$ or $v=-\alpha$, the original equation has the root $0$
or $2u$, respectively. Horizontal lines meet each member by bijectivity
of cubing. Vertical lines and the line at infinity contain $[0:1:0]$.
Consequently $N(S)$ is the exact number of missing projective lines for
the cusp. Translating $t=x+u$ in the same equation also shows that every
blocking set of parameters for the cubic graphs blocks the cusps.

\subsection{The two meridians at each puncture}

The product of the ordering torsors of the three roots is finite
\'etale on $U$, with group $G=S_3^\ell$. We compute its geometric
monodromy first over $\mathbb C$ with arbitrary distinct finite marks
$a_i=-\alpha_i$. Add the mark $a_\infty=\infty$ on the sphere. Place
the marks on a topological equator and place the base positions of
$v,w$ in opposite hemispheres. Join each base position to the marks
by a fan of arcs in its hemisphere. Moving $v$ round one mark using
its arc gives $x_i$ and moving $w$ gives $y_i$, with indices including
$\infty$. For $i\ne j$ the two loops have disjoint tubes, so that
$x_i y_j=y_j x_i$. These loops generate $\pi_1(U)$. Indeed, the fibre
of the map forgetting $v$ is a sphere punctured at all marks and at
$w$. Its meridians around the marks generate, by the boundary relation.
The $y_i$ project to generators of the base. The homotopy sequence of
the configuration-space fibration gives the assertion \cite{FN}.

The cubic cover $T\mapsto T^3-3T$ has transposition inertia at
$-2,2$ and a three-cycle at infinity. For fixed $w$, the map
$v\mapsto t_{\alpha_i}$ is fractional linear, carrying $a_i,\infty,w$
to $-2,2,\infty$. Thus a finite $x_i$ acts by a transposition $A_i$
in coordinate $i$ alone. The analogous transposition for $y_i$ is
$B_i$. In coordinate $i$, $x_\infty$ is a transposition commuting with
$B_i$, hence equals $B_i$. The monodromy on the fixed-$w$ fibre is
transitive and contains a transposition, so equals $S_3$; its image is
generated by $A_i$ and the image of $x_\infty$. Thus $A_i\ne B_i$.
The finite meridians generate every coordinate separately, giving
geometric monodromy $G$.

\subsection{The degree-one cohomology}

Let $W$ run over the external tensors obtained by taking the trivial,
sign or standard representation in each coordinate, over a field of
characteristic zero. Write $b(W)=\dim H^1(U(\mathbb C),W)$. We claim
\begin{equation}\label{eq:cubic-betti}
 \sum_{W\ne{\bf1}} b(W)\le 10\ell\,2^\ell .
\end{equation}
For a crossed homomorphism $D$, put $u_i=D(x_i)$, $v_i=D(y_i)$.
The commuting pairs impose
\begin{equation}\label{eq:cubic-meridians}
 (1-B_j)u_i=(1-A_i)v_j\qquad (i\ne j),
\end{equation}
where for the infinite index the letters denote the full meridian
actions. Coboundaries have $D(g)=(1-g)z$. These necessary relations
on a generating set suffice for an upper bound.

First suppose the tensor has no sign and at least one standard factor.
For finite $i$, the projection of $u_i$ to $W/(1-A_i)W$ is invariant
under every $B_j$ with finite $j\ne i$. The quotient in factor $i$ and
the fixed spaces in the other factors all have dimension at most one.
The analogous statement for $v_i$ gives at most $2\ell$ dimensions in
all, and these projections vanish on coboundaries. In their kernel
write $u_i=(1-A_i)P_i$, $v_j=(1-B_j)Q_j$. In each standard factor use
as basis $e,f$ on the distinct fixed lines of $A_i,B_i$, respectively.
The quotient by $\ker(1-A_i)$ retains only words with $f$ in position
$i$; the quotient for $B_j$ retains only words with $e$ in position $j$.
Trivial factors retain no coefficient at their own index. For $i\ne j$
the two operators act on different factors, so the kernel of their
product is the sum of their kernels. Equation~\eqref{eq:cubic-meridians}
therefore identifies the coefficients of $P_i,Q_j$ on every word
retained by both quotients.

For a mixed word the indices occupied by $f$ and by $e$ form two
nonempty disjoint sets. Every left index is joined to every right one,
so all retained coefficients of that word agree. Subtract a single
coboundary whose $z$ has these common coefficients on all mixed words.
The two constant words leave at most $2\ell$ coefficients altogether.
With the finite-loop data fixed, the difference of two choices for
$u_\infty$ is fixed by every finite $B_j$, a space of dimension at most
one. The same reasoning applies to $v_\infty$. This gives $b(W)\le4\ell+2$.

If there is a sign at $s$, subtract the unique coboundary making
$u_s=0$ and put $h=v_s$. For $j\ne s$, including $j=\infty$,
\[
          v_j=0,\qquad 2u_j=(1-A_j)h.
\]
Thus a cohomology class is determined by $h$. For distinct finite
$i,j\ne s$ we have $(1-A_i)(1-B_j)h=0$. With exactly one sign this
excludes every mixed word on the standard factors, giving $b(W)\le2$.
With another sign at $t$, the equations with $(i,t)$ and $(t,i)$ force
any standard factor $i$ to lie in both its fixed lines, so $h=0$.
Two signs with all remaining factors trivial leave at most one
dimension. Three or more signs give $4h=0$ by taking two of them
different from $s$. The resulting sum is at most
\[
 (4\ell+2)(2^\ell-1)+2\ell\,2^{\ell-1}+\binom{\ell}{2}
 \le 10\ell\,2^\ell ,
\]
which proves \eqref{eq:cubic-betti}.

\subsection{Specialization and the trace estimate}

Lift the marks to disjoint sections over the localization of
$\mathbb Z[\zeta_{q-1}]$ at a prime above $p$ using roots of unity and
the constant sections $0,\infty$. The residue field is $\mathbb F_q$:
if $q=p^h$, the order of $p$ modulo $p^h-1$ is $h$. In
$\mathbb P^1\times\mathbb P^1$ the removed verticals, horizontals and
diagonal have triple intersections precisely at $(a,a)$ for the marks
$a$, including infinity. Blow up those disjoint sections. The three
tangents are distinct in both fibres, giving a smooth proper relative
surface with relative simple normal crossings boundary. On its open
complement the ordered-root torsor extends because each discriminant
is a unit. Its order $6^\ell$ is prime to $p$, so the local systems it
trivializes are tame, including at the exceptional divisors.

Over the strictly henselian base, apply the tame base-change theorem
in degrees zero and one \cite[Exp.~XIII, Cor.~2.9]{SGA1}. First apply
it to the sheaf of $\mathbb F_2$-valued functions on the torsor. The
dimension of $H^0$ counts geometric connected components, so the
complex computation gives full geometric monodromy in the special
fibre as well. Apply the same theorem to successive quotients modulo
$2^j$ of an invariant $\mathbb Z_2$-lattice in each $W$. The finite
degree-zero groups satisfy the Mittag--Leffler condition. Passing to
the inverse limit and inverting $2$, then using \'etale--singular
comparison in characteristic zero, transports \eqref{eq:cubic-betti}
to $\overline{\mathbb F}_q$.

For nontrivial $W$ of rank $r$, affine vanishing, Poincar\'e duality
and the absence of geometric invariants give $H_c^0=H_c^1=H_c^4=0$
and $\dim H_c^3=b(W)$; the representations are self-dual. A nonidentity
deck transformation of the $G$-torsor acts freely and has order prime
to $p$. Its alternating compactly supported trace vanishes by the
finite-order trace theorem \cite[Thm.~3.2, Prop.~3.12]{DL}. The
equivariant Euler characteristic is consequently a multiple of the
regular character. Taking invariants identifies the multiple as
$\chi_c(U)=(1-\ell)^2-(1-\ell)=\ell(\ell-1)$, by subtracting the
diagonal. It follows that $\dim H_c^2=r\ell(\ell-1)+b(W)$.
Finite monodromy gives weight zero. The trace formula and the weight
bound \cite[\S3.3]{WeilII} give
\[
 \left|\sum_{z\in U(\mathbb F_q)}
             \operatorname{tr}(\operatorname{Frob}_z\mid W)\right|
 \le r\ell(\ell-1)q+b(W)(q+q^{3/2}).
\]
On $S_3$, three times the indicator of the three-cycles is the character
$1+\mathrm{sgn}-\mathrm{std}$. Expand the product over the $\ell$
coordinates. The trivial term contributes $(q-\ell)(q-\ell-1)$, and
the sum of the nontrivial ranks is $4^\ell-1$. The preceding estimate
and \eqref{eq:cubic-betti} give \eqref{eq:cubic-error}.

After division by $q^2$ the error is at most
\[
       \frac{\ell^2 4^\ell}{q}+
       \frac{10\ell\,2^\ell}{\sqrt q}(1+q^{-1/2}).
\]
It tends to zero uniformly for $1\le\ell\le cL$ if $c<1/\log4$.
At $\ell=\lfloor(L-2\log L)/\log4-5\rfloor$, the upper limit of
this error is at most $4^{-5}/(\log4)^2+10\cdot2^{-5}/\log4<0.23$.
The normalized main term tends to one. Thus every set of that size has
a missing line, proving the strict lower bound in
Theorem~\ref{thm:cubic-surface}.

For the upper bound, $x^3-mx$ omits exactly $(q+1)/3$ values for
$m\ne0$. Indeed, $x^2+xy+y^2=m$ is a norm conic with $q+1$ ordered
solutions, since the nontrivial cube roots of unity are quadratic over
$\mathbb F_q$. Let $h$ count its diagonal solutions. The cubic has $h$
fibres of size two, one at each rational critical value. If $k$ fibres
have size three, counting ordered unequal collisions gives
$2h+6k=q+1-h$. Hence the number of empty fibres is $h+2k=(q+1)/3$.
Horizontal lines are all hit. With $j$ independent uniform translates
the expected number of missed lines is at most
$q(q-1)((q+1)/(3q))^j$, which yields the claimed upper limit.
The transfer to the cusp was proved above.

The fixed-length tuple asymptotic \cite{BaryTuple} and the exact cubic
twin-prime count \cite{BaryStix} concern related irreducibility
conditions. The estimate needed here controls every configuration of
$\ell$ marks with $\ell$ growing. The gain in \eqref{eq:cubic-error}
comes from summing the first Betti numbers separately from the ranks.

For orientation, applying a one-variable trace estimate to a slope with
disjoint critical values gives the smaller coefficient $1/\log16$:
there are $2\ell$ finite branch points, the nontrivial tensor ranks
sum to $4^\ell-1$, and the tame curve estimate gives an error at most
$(2\ell-1)(4^\ell-1)\sqrt q$ in the normalized count multiplied by
$3^\ell$. This follows from the character estimate in
\cite[Prop.~5.1]{Zywina} and the full-cycle norm in
\cite[Props.~6,10]{Bellaiche}. More general explicit Chebotarev
estimates are also available \cite{Shpilka}; the multishift
specialization in its Theorem~1.6 assumes degree at least four.

\subsection{Earlier tuple estimates and two further cubic bounds}

For $p>3$, Pollack's estimate for an individual cubic already gives a
logarithmic lower bound. Apply \cite[Lem.~6, Cor.~15 and the estimate
preceding (4.4)]{Pollack} to the outer polynomials $T-\alpha$. At least
$q^2-36q(1+\binom{\ell}{2})$ cubics $h=T^3+aT^2+eT$ are admissible.
For each, at least
$3^{-\ell}[q-2(g\sqrt q+g+6^\ell)]-2\ell$ values of $v$ make all
$h-v-\alpha$ irreducible, with $g\le3\ell6^\ell$. Translation by
$-a/3$ removes the quadratic term. This proves the lower coefficient
$1/(2\log6)$ directly from Pollack, even with $\ell$ growing.
The two-parameter estimate \eqref{eq:cubic-error} gives $1/\log4$.

In characteristic two take $q=2^{2j+1}$. At slope one the substitution
$x=z+z^{-1}$ gives $x^3+x=z^3+z^{-3}$. Inversion and multiplication
by cube roots of unity give its $S_3$-closure. At infinity inertia is
tame of order three. At the single finite branch value, inversion at
$z=1$ changes $z-1$ by a term of valuation two, hence
$I_0=I_1=C_2$ and $I_2=1$. The $\ell$ translated finite branch values
are distinct, so the intersection argument of Section~\ref{sec:fixed-graph}
gives geometric group $S_3^\ell$. For a nontrivial tensor of rank $r$
the total Swan conductor is at most $\ell r$. The Euler characteristic
formula on the punctured affine line consequently gives
$\dim H_c^1\le(2\ell-1)r$. If $N_1$ counts intercepts missing every
translate, the same character expansion yields
\[
  |3^\ell N_1-(q-\ell)|\le(2\ell-1)(4^\ell-1)\sqrt q .
\]
This gives $\liminf B(q,3)/\log q\ge1/\log16$ along these fields.

For all $q\equiv2\pmod3$, the graph upper bound refines to
$B(q,3)\le(2L-\log L)/\log3+17$. In characteristic two the norm conic
still has $q+1$ points and one diagonal point; the fibre at zero is
the unique two-point fibre. The same single-line miss rate follows.
After $k$ independent samples, at most $t$ lines can remain uncovered
for some sample whenever $q(q-1)((q+1)/(3q))^k\le t$. One further
translate hits all $t$ lines if
$q+1>[2+(2t-1)3^t]\sqrt q+3^t$. Indeed, put $\alpha=z^{-3}$ and
lift the root $1$ of $Y_i^3-1-m_iY_iz^2-b_iz^3$ by Hensel's lemma.
The roots $Y_i/z$ generate over $\mathbb F_q(\alpha)$ a function field of degree at most
$3^t$ with full constants in $\mathbb F_q((z))$. Successive
Castelnuovo--Severi bounds give twice its genus at most
$2+(2t-1)3^t$. Hasse--Weil then gives a rational place away from
the at most $3^t$ poles of $\alpha$. For $L\ge128$ take
$t=\lfloor L/(4\log3)\rfloor$ and $k=\lfloor(2L-\log L)/\log3\rfloor+16$.
Here $t\ge L/16$, $k/q<1$, $3^t\le q^{1/4}$ and
$2t+4\le q^{1/4}$, which imply both inequalities.
For $q\ge16$ with $L<128$, $k=\lceil2L/\log3\rceil+3$ samples already
give expected misses below one, since $k/q<1$. For $q<16$ use all
parameters. The companion manuscript, \S\S12--13, gives the detailed
genus and rational-place bounds.

\section{Different blocking numbers on the same field}\label{sec:norm-spheres}

Let $F=\mathbb F_s\subset E=\mathbb F_{s^n}$, and let $N=N_{E/F}$.
Write $B_{\rm cusp}(q,d)$ for the pencil $X^d=(Y+\alpha Z)^{d-1}Z$
and put $B_{\rm graph}(q,d)=B(q,d)$ for the pencil $YZ^{d-1}=X^d-\alpha Z^d$.

\begin{theorem}\label{thm:sphere-minimum}
In relative degree two, for every prime power $s\ge2$,
\[
   B_{\rm cusp}(s^2,s)=3,\qquad B_{\rm graph}(s^2,s)=2s-1 .
\]
In relative degree three, for $s\ge4$,
\[
   B_{\rm cusp}(s^3,s)=4,\qquad B_{\rm graph}(s^3,s)=3s-2 .
\]
Both minima in each pair are realized by nontrivial projective blocking
sets. In particular the ratio of the two minima is unbounded.
\end{theorem}

Call $S\subset E$ cospherical if there exist $b\in E$ and
$R\in F^\times$ such that $N(b+\alpha)=R$ for every $\alpha\in S$.
The following reduction also holds outside the two degrees of the theorem.

\begin{lemma}\label{lem:cospherical}
Every set that fails to block the cusp of degree $s$ is cospherical.
If $|S|<s$, a cospherical set fails to block.
\end{lemma}
\begin{proof}
Assume $S\ne\varnothing$, the empty case being immediate. A missing
line $y=mx+b$ has $m\ne0$ because $t\mapsto t^s$ is onto, and all
$z=b+\alpha$ are nonzero because $t=0$ gives the value zero. For
$v=t^{-1}$ the intercept equation becomes
\[
                     zv^s+mv=1.
\]
Every solution has $v\ne0$. This $F$-linear operator is invertible
unless $N(z)=N(-m)$: a nonzero kernel vector exists precisely when
$-m/z$ is an $(s-1)$st power, equivalently when $N(-m/z)=1$.
Thus every omitted line forces a common nonzero norm.

Conversely, suppose $N(z)=R$ for all $z=b+\alpha$ and choose one such
element $a$. Each map $v\mapsto zv^s-av$ has an $F$-line as its
kernel, hence image of size $q/s$. If $|S|<s$, some $u\in E$ belongs to
none of the images, and $u\ne0$. Take $k\in E$ with $k^s=u$ and put
$m=-ak/u$. A solution of $zv^s+mv=1$ would give
$z(kv)^s-a(kv)=u$. This excludes every shift at once.
\end{proof}

For any fixed $b$, the polynomial $N(b+T)=\prod_{j=0}^{n-1}
(T+b^{s^j})$ in $F[T]$ is monic of degree $n$. Any $n+1$ distinct
elements of $F$ therefore block by Lemma~\ref{lem:cospherical}.
Conversely, in degree at least two, any two distinct elements $x,y$
are cospherical. Choose $w\ne1$ of norm one, put $v=(y-x)/(w-1)$,
and translate by $v-x$. The two values $v,wv$ have the same nonzero norm.
This proves the quadratic cusp minimum for $s\ge3$, including $s=3$,
where the upper bound uses the first assertion of the lemma. When
$s=2$, both conic pencils have exactly the same homogeneous equation;
\eqref{eq:graph-frob} gives the common value $B(4,2)=3$.

In degree three we show that every three distinct points are
cospherical. By translation we may take them to be $0,a,c$; put
$u=c/(a-c)$, so $u(1+a/z)=(1+u)(1+c/z)-1$. For any $A,B\in
F^\times$, consider
\[
             g_j(T)=T^3+jT^2+(j+B-A-1)T-A,\qquad j\in F .
\]
Their values at $0,-1$ are $-A,-B$. Every other point of $F$ excludes
precisely one $j$, as its coefficient in the root equation is
$T(T+1)\ne0$. Thus at least two $g_j$ have no root in $F$ and are
irreducible. Each splits in $E$. Any root $y$ of such a cubic has
$N(y)=A$, $N(1+y)=B$. Take $A=N(u)$ and $B=N(1+u)$, and choose $y\ne
u$; the two irreducibles have disjoint root sets. Set $z=ua/(y-u)$.
Then $N(z+a)=N(z)=N(z+c)\ne0$ by the two displayed identities.
The lemma proves the cubic lower bound. This argument uses only the
degree-three case of simultaneous norm equations.
The existence of such sphere centres already follows, for all $s$, from
the common-neighbour formulas of Bayer, M\'esz\'aros, R\'onyai and
Szab\'o for the projective norm graph $\mathrm{NG}(s,4)$
\cite[Thm.~1(a), Lem.~2(iv), Prop.~1]{Bayer}. In their recurrence
the quadratic two-norm fibres have at most two points, giving at least
$2(s+1)-2(s-2)=6$ solutions before exclusion of one point. Their
vertices $(\alpha_i,1)$ therefore have at least five common neighbours
when the three first coordinates are distinct. A neighbour $(b,R)$
specifies precisely such a centre and radius. We included the cubic
argument to give an elementary kernel-verified proof. The
blocking-number application also requires the simultaneous missed-line
criterion and a blocking set attaining the upper bound.

For comparison, the Frobenius graph formula gives
$B_{\rm graph}(s^n,s)=n(s-1)+1$. In both examples the graph minimum
is $<q$, so its optimal unions contain no full rational line by the
projective graph criterion. The cusps all have the single point
$[0:1:0]$ at infinity. At each fixed $x$, the equation
$(y+\alpha)^{s-1}=x^s$ admits at most $s-1$ values of $y$ per shift.
For the stated sizes, $3(s-1)<s^2$ and $4(s-1)<s^3$, so even the
vertical lines cannot be contained in the union. All other projective
lines have a point at infinity outside the union.

\section*{Proof verification and reproducibility}
The accompanying Lean~4.8 project uses the homogeneous projective
equations of the two pencils. It proves the two exact comparisons of
Theorem~\ref{thm:sphere-minimum} over the same extension field, the
projective nontriviality bounds, and the cubic simultaneous-norm
construction, without geometric or analytic hypotheses. It checks
geometric irreducibility over the algebraic closure for both pencils. The
general cuspidal parametrization, the equivalence with the slope test,
and the cuspidal logarithmic upper bound are also unconditional.

For the cuspidal lower bound the kernel checks off-coordinate
fixedness for the inertia subgroup of the actual product splitting
field, the local Eisenstein factor, Hensel lifting of the other
sheets, and restriction of the resulting permutations to that field.
The proof of the uniform count uses one threshold for all shift sets
in the logarithmic range. Lean derives the degree and genus budgets
from the individual root steps. Applying Castelnuovo--Severi to the
actual function fields, constructing the arithmetic twist, and
applying Hasse--Weil and Frobenius specialization remain geometric
inputs. The wild graph argument also uses the geometric group
identifications of Section~\ref{sec:fixed-graph}.
For the cubic surface estimate, Lean checks the actual pair-count
reduction, the linear algebra of crossed cocycles under the stated
meridian relations, and the passage from the cohomological input to
the $1/\log4$ bound. The identification of the meridians, specialization
of the marked surface, the trace formula and the weight estimates are
geometric inputs. The cubic upper bound with the negative $\log\log q$
term is unconditional in Lean for $\log q\le3^9$. Its all-field
formalization uses the genus and point theorems on the constructed
Laurent branch. The \href{https://github.com/Heinz1729/monomial-blocking-curves}{public
repository} includes the sources, finite-field checks, proof-scope
index and a reproducible audit of theorem dependencies.

\section*{Attribution}
GPT-6.2 Astra is the AI byline of this public preprint. The research
chat supplied the proposals for the cuspidal pencil and the norm-sphere
comparison. OpenAI's Codex was used for proof development, Lean
formalization, source checks and preparation of the manuscript.

\begin{thebibliography}{99}
\raggedright

\bibitem{AGY} S. Asgarli, D. Ghioca, C.~H. Yip, Blocking sets from a
union of plane curves, Des. Codes Cryptogr. 94 (2026), article 217.
\href{https://doi.org/10.1007/s10623-026-01958-z}{doi:10.1007/s10623-026-01958-z}.
\bibitem{Jamison} R.~E. Jamison, Covering finite fields with cosets of
subspaces, J. Combin. Theory Ser.~A 22 (1977), 253--266.
\href{https://doi.org/10.1016/0097-3165(77)90001-2}{doi:10.1016/0097-3165(77)90001-2}.
\bibitem{BS} A.~E. Brouwer, A. Schrijver, The blocking number of an
affine space, J. Combin. Theory Ser.~A 24 (1978), 251--253.
\href{https://doi.org/10.1016/0097-3165(78)90013-4}{doi:10.1016/0097-3165(78)90013-4}.
\bibitem{Kurlberg} P. Kurlberg, Poisson spacing statistics for value
sets of polynomials, Int. J. Number Theory 5 (2009), 489--513.
\href{https://arxiv.org/abs/math/0602673}{arXiv:math/0602673}.
\bibitem{Pollack} P. Pollack, Simultaneous prime specializations of
polynomials over finite fields, Proc. Lond. Math. Soc. (3) 97 (2008),
545--567. \href{https://doi.org/10.1112/plms/pdn013}{doi:10.1112/plms/pdn013}.
\bibitem{Reid} C.~D. Reid, Multiple transitivity except for a system of
imprimitivity, J. Group Theory 27 (2024), 651--712.
\href{https://doi.org/10.1515/jgth-2023-0062}{doi:10.1515/jgth-2023-0062}.

\bibitem{FN} E. Fadell, L. Neuwirth, Configuration spaces, Math. Scand.
10 (1962), 111--118.
\href{https://doi.org/10.7146/math.scand.a-10517}{doi:10.7146/math.scand.a-10517}.
\bibitem{SGA1} A. Grothendieck et al., Rev\^etements \'etales et groupe
fondamental (SGA 1), Expos\'e XIII, Lecture Notes in Mathematics 224,
Springer, 1971. \href{https://arxiv.org/abs/math/0206203}{Updated edition}.
\bibitem{DL} P. Deligne, G. Lusztig, Representations of reductive groups
over finite fields, Ann. of Math. (2) 103 (1976), 103--161.
\href{https://doi.org/10.2307/1971021}{doi:10.2307/1971021}.
\bibitem{WeilII} P. Deligne, La conjecture de Weil. II, Publ. Math.
IH\'ES 52 (1980), 137--252.
\href{https://doi.org/10.1007/BF02684780}{doi:10.1007/BF02684780}.
\bibitem{BaryTuple} L. Bary-Soroker, Hardy--Littlewood tuple conjecture
over large finite fields, Int. Math. Res. Not. IMRN 2014 (2014), 568--575.
\href{https://doi.org/10.1093/imrn/rns249}{doi:10.1093/imrn/rns249}.
\bibitem{BaryStix} L. Bary-Soroker, J. Stix, Cubic twin prime polynomials
are counted by a modular form, Canad. J. Math. 71 (2019), 1323--1350.
\href{https://doi.org/10.4153/CJM-2018-018-9}{doi:10.4153/CJM-2018-018-9}.

\bibitem{ZhouWu} J. Zhou, H. Wu, Covering finite fields by translates of
multiplicative subgroups, preprint (2026).
\href{https://arxiv.org/abs/2610.25228}{arXiv:2610.25228}.
\bibitem{Stichtenoth} H. Stichtenoth, Algebraic Function Fields and Codes,
second edition, Graduate Texts in Mathematics 254, Springer, 2009.
\href{https://doi.org/10.1007/978-3-540-76878-4}{doi:10.1007/978-3-540-76878-4}.
\bibitem{Kosters} M. Kosters, A short proof of the Chebotarev density
theorem for function fields, Math. Commun. 22 (2017), 227--233.
\href{https://hrcak.srce.hr/file/274189}{Journal article}.
\bibitem{Stacks} The Stacks Project Authors, Galois extensions and
ramification, \href{https://stacks.math.columbia.edu/tag/09E3}{Tag 09E3}
and \href{https://stacks.math.columbia.edu/tag/09EC}{Tag 09EC}.
\bibitem{Bellaiche} J. Bella\"iche, Th\'eor\`eme de Chebotarev et
complexit\'e de Littlewood, Ann. Sci. \'Ec. Norm. Sup\'er. (4) 49
(2016), 579--632.
\href{https://doi.org/10.24033/asens.2291}{doi:10.24033/asens.2291}.
\bibitem{Zywina} D. Zywina, Hilbert's irreducibility theorem and the
larger sieve, preprint (2010).
\href{https://arxiv.org/abs/1011.6465}{arXiv:1011.6465}.
\bibitem{Shpilka} I. Benbenishty Shpilka, Explicit error term in the
Chebotarev density theorem for function fields, preprint (2026).
\href{https://arxiv.org/abs/2611.28155}{arXiv:2611.28155}.
\bibitem{Bayer} T. Bayer, T. M\'esz\'aros, L. R\'onyai, T. Szab\'o,
Subgraphs of the projective norm graph, Forum Math. Sigma 14 (2026), e121.
\href{https://doi.org/10.1017/fms.2026.10245}{doi:10.1017/fms.2026.10245}.
\end{thebibliography}
\end{document}
